% sections/04-frontend-interface.tex — Frontend Interface

\section{Frontend Interface}

% ── Requirement ──
\subsection{Requirement}

Your project must implement at least \textbf{one} user interface. The interface should allow users to interact with your application and visualize the results.

\begin{notebox}[Bonus Points]
Implementing \textbf{two or more} interface types earns \textbf{+2 bonus points}.
\end{notebox}

% ── Interface Options ──
\subsection{Interface Options}

\subsubsection{\cpp{} GUI}

\begin{multicols}{2}
\begin{itemize}
  \item \textbf{Qt} --- Full-featured, cross-platform GUI framework with rich widget library and designer tools.
  \item \textbf{ImGui} --- Immediate-mode GUI, ideal for debug/visualization tools with minimal setup.
  \item \textbf{wxWidgets} --- Native look-and-feel, cross-platform widget toolkit.
  \item \textbf{FLTK} --- Lightweight, fast, small footprint GUI library.
\end{itemize}
\end{multicols}

\subsubsection{Python Frontend}

Use Python for the interface while keeping the core logic in \cpp{}. This approach requires a binding layer between the two languages (see Section~\ref{sec:integration-tools}).

\subsubsection{Web Interface}

Build a \cpp{} backend serving a REST API with an HTML/CSS/JavaScript frontend. Recommended frameworks:

\begin{itemize}
  \item \textbf{Crow} --- Micro web framework for \cpp{}, Flask-inspired.
  \item \textbf{Drogon} --- High-performance async web framework.
  \item \textbf{cpp-httplib} --- Single-header HTTP/HTTPS library.
\end{itemize}

\subsubsection{Desktop Application}

\begin{itemize}
  \item \textbf{Electron} --- Web technologies (HTML/CSS/JS) with \cpp{} native addon via N-API.
  \item \textbf{Tauri} --- Rust-based, lighter alternative to Electron with lower resource usage.
\end{itemize}

% ── C++/Python Integration Tools ──
\subsection{\cpp{}/Python Integration Tools}
\label{sec:integration-tools}

\subsubsection{ctypes}

Python's built-in foreign function library for calling C-compatible shared libraries.

\begin{lstlisting}[style=python]
import ctypes

# Load the shared library
lib = ctypes.CDLL("./libproject.so")

# Define argument and return types
lib.search.argtypes = [ctypes.c_char_p]
lib.search.restype = ctypes.c_int

# Call the function
result = lib.search(b"query")
print(f"Result: {result}")
\end{lstlisting}

Best for: Simple C-style interfaces, quick prototyping.

\subsubsection{cffi}

C Foreign Function Interface for Python, providing a cleaner API than ctypes.

\begin{lstlisting}[style=python]
from cffi import FFI

ffi = FFI()
ffi.cdef("""
    int search(const char* query);
    void insert(const char* key, int value);
""")

lib = ffi.dlopen("./libproject.so")
result = lib.search(b"query")
\end{lstlisting}

Best for: C libraries, cleaner API than ctypes.

\subsubsection{pybind11}

Seamless \cpp{}/Python interoperability with modern \cpp{} features.

\begin{lstlisting}[style=cpp]
#include <pybind11/pybind11.h>
namespace py = pybind11;

class HashTable {
public:
    void insert(const std::string& key, int value);
    int search(const std::string& key);
};

PYBIND11_MODULE(project, m) {
    py::class_<HashTable>(m, "HashTable")
        .def(py::init<>())
        .def("insert", &HashTable::insert)
        .def("search", &HashTable::search);
}
\end{lstlisting}

Best for: \cpp{} classes, STL containers, modern \cpp{}.

\subsubsection{SWIG}

Simplified Wrapper and Interface Generator. Generates bindings for multiple target languages from annotated \cpp{} headers.

Best for: Multi-language support, large existing codebases.

\subsubsection{Boost.Python}

Part of the Boost libraries, provides a rich interface for \cpp{}/Python interoperability.

Best for: Projects already using Boost libraries.

% ── Recommendations ──
\subsection{Recommendations}

\begin{enumerate}
  \item \textbf{Simple projects:} \cpp{} Console or ctypes
  \item \textbf{Visualization-heavy:} Qt or ImGui
  \item \textbf{Python integration:} pybind11
  \item \textbf{Web deployment:} REST API + Web frontend
\end{enumerate}
